\section{Coming into the Light in Confession}

In the sacrament of Penance we come before Christ's merciful judgment while conversion is still possible. We examine conscience, confess our sins, receive absolution, and undertake the penance given. The Catechism calls the sacrament, in a certain way, an anticipation of the judgment at the end of life (CCC 1470). What we fear acknowledging can be brought to the Savior now. He has instituted this means of forgiveness for sinners who turn to him.

Examination begins by asking for truth before asking for a favorable account. The commandments, the Gospel, and the duties of our state of life give conscience something firmer than the impression we happen to have of ourselves. We consider thoughts, words, deeds, and omissions, the wrong we chose and the good we owed. We also remember the grace which accompanied us. Gratitude helps us recognize where God has acted; it prevents examination from becoming a story in which sin seems more powerful than the mercy that has already sustained us.

A useful examination moves from general characterizations to particular acts. It is easy to say that we have been selfish and difficult to acknowledge the money withheld, the responsibility abandoned, or the harsh word deliberately chosen. General remorse can leave the protected fault intact. Naming the act allows us to see what repentance requires. The purpose is neither a polished narrative nor a complete psychological autobiography. We bring our sins honestly, with the circumstances necessary to understand them, and receive the priest's help where we are uncertain.

Francis de Sales's meditation on judgment concludes with examination and confession. After contemplating the separation of the good and evil, he considers the opening of conscience: the malice of sin is disclosed, but so are repentance and the effects of grace. His practical resolution is to examine conscience, confess faults, and follow the counsel received (\emph{Introduction to the Devout Life} I.14, considerations~5--6 and affections~2). Someone following that counsel may discover a fault he had avoided naming. The time spent meditating has then become preparation for an act of repentance.

Contrition is sorrow for sin with the resolution to turn from it. It can begin in fear of judgment and punishment; such imperfect contrition is itself a prompting of grace and disposes us for sacramental forgiveness. Perfect contrition arises from love of God above all things and includes the firm resolve to seek sacramental confession as soon as possible (CCC 1451--1453). We need not reject the first movement of fear because our love is still weak. We can let it bring us to Christ, asking that sorrow deepen into grief for having offended the One who loves us.

The sincerity of this sorrow is not measured by tears or by the intensity of emotion. We may feel profoundly and remain unwilling to change; we may feel little and make a costly, deliberate act of repentance. The will's return matters. A firm purpose of amendment means that we intend to reject the sin and take the steps required to resist it. It does not mean that we can guarantee, by our own strength, that no weakness will ever reappear. We ask for grace and choose the means of fidelity which are available to us.

Amendment often requires changing circumstances in which we repeatedly consented to sin. If a conversation regularly becomes detraction, we can decide beforehand to stop participating when it does. If a duty is repeatedly neglected, we may need to set aside a definite time for it and relinquish something which displaces it. The intention to change should include willingness to accept such consequences. A confessor's counsel can help make that intention practical and proportionate.

Confession requires honesty about grave sin, including what was secret. The shame that tempts us to conceal a fault can be offered as part of the return. Yet sincere confession is not an achievement of perfect memory. We examine diligently and confess what we know, without deliberately withholding what should be brought before God's mercy (CCC 1454--1456). Forgetfulness and concealment are different. The penitent's task is to be truthful and receptive, accepting guidance when a doubt needs resolution rather than attempting to manufacture certainty by endless recollection.

Absolution restores communion with God and reconciles the sinner with the Church (CCC 1468--1469). We should receive this as an action of Christ. The priest's ministry brings an objective gift which does not wait upon a corresponding feeling of relief. A person may leave confession still emotionally troubled and nevertheless truly reconciled. Scrupulous distress calls for patient pastoral care and, where appropriate, professional help. Its insistence is not proof that forgiveness failed. Fidelity can include the humble acceptance of sound counsel when private anxiety demands that an honestly confessed matter be reopened without end.

Forgiveness also leads toward satisfaction and reparation. Stolen goods should be returned, harm to a reputation repaired where possible, and injuries addressed according to justice. Absolution removes sin; the disorder caused by sin may still require healing (CCC 1459--1460). The penance assigned is part of that return, and the wider duty to repair harm belongs to love of the neighbor. We should not seek a sacramental peace which depends upon leaving someone else to bear an avoidable loss caused by our wrongdoing.

Reparation can require prudence. An indiscriminate disclosure may cause further injury. Some restitution is straightforward; other cases require confidential counsel about what is owed and how to fulfill it. Difficulty should lead us to seek that guidance. If shame makes us reluctant, we can begin by acknowledging the reluctance and asking for help. Forgiveness restores us to the freedom to do what justice and charity require, including a repair whose cost we would prefer to avoid.

Returning to confession also teaches perseverance after a fall. The recurrence of a temptation can expose a weakness in the means we chose or a reluctance to use them. We can bring that failure honestly, accept counsel, and begin again. A penitent need not wait until he can present a history of unbroken success. He must come with sincere sorrow and a purpose of amendment, trusting the forgiveness Christ gives through his Church.

